\section{Estrategia de Validación y Testing}

Para asegurar que el modelo de la bomba de infusión cumple rigurosamente con los requisitos especificados, se implementó una estrategia de validación integral dividida en múltiples capas. Esta estrategia garantiza tanto el correcto funcionamiento individual de las partes como el comportamiento global del sistema acoplado.

\subsection{Pruebas Unitarias (Unit Tests)}
Las pruebas unitarias tienen como objetivo verificar la lógica interna de cada modelo atómico DEVS de forma aislada. Mediante la inyección de secuencias de entrada controladas directamente en los puertos de cada componente (como el \textit{Controlador}, \textit{Sensor de Flujo}, \textit{Actuador}, etc.), se validó que las funciones de transición (interna y externa) y las funciones de salida generen los estados correctos. Esto asegura que la base del sistema esté libre de errores lógicos antes de su integración.

\subsection{Pruebas de Integración y Construcción de Escenarios}
Para evaluar el modelo acoplado completo, se diseñó una herramienta denominada \texttt{ScenarioRunner}. Esta clase actúa como un orquestador que permite configurar y ejecutar simulaciones determinísticas sin alterar el código fuente del modelo formal. 

Los escenarios se construyen mediante la técnica de \textit{monkey patching} controlada temporalmente, lo que permite:
\begin{itemize}
    \item \textbf{Inyectar órdenes médicas:} Programar cambios de caudal en instantes de tiempo específicos.
    \item \textbf{Simular fallas físicas:} Alterar la salida del sensor de flujo para emular discrepancias entre el caudal real y el medido.
    \item \textbf{Controlar el entorno humano:} Silenciar o programar las confirmaciones del personal de enfermería para evaluar la escalabilidad de las alarmas.
\end{itemize}
De esta forma, el sistema se somete a condiciones de estrés predefinidas que emulan situaciones clínicas reales.

\subsection{Verificación Automática de Propiedades}
En lugar de verificar los resultados visualmente, se desarrolló un módulo de \textit{Property Checkers} (Verificadores de Propiedades). Este módulo analiza las trazas temporales extraídas por un monitor pasivo y evalúa el cumplimiento de las reglas del proyecto de manera automática.

Es fundamental destacar que \textbf{estos verificadores son genéricos}. Al estar desacoplados de la lógica de inyección de los escenarios, pueden recibir una traza generada por un escenario estrictamente controlado, o bien, una traza resultante de una simulación estocástica prolongada (con tiempos y fallas aleatorias), y verificar con la misma rigurosidad si el sistema violó alguna regla.

Las propiedades verificadas se dividen en tres categorías principales basadas en la especificación formal del proyecto:
\begin{itemize}
    \item \textbf{Propiedades de Seguridad (Safety):} Aseguran que el sistema no cause daño (e.g., el caudal no supera los 200 ml/h, y la bomba no infunde líquido bajo una orden de detención o durante un bloqueo crítico).
    \item \textbf{Propiedades de Vivacidad (Liveness):} Garantizan que el sistema eventualmente tome acciones esperadas (e.g., una alarma crítica no confirmada repite su notificación de forma indefinida).
    \item \textbf{Propiedades Temporales:} Verifican el cumplimiento de restricciones de tiempo estricto (e.g., inicio de infusión en $<3$ segundos, detención a los 60 segundos tras el fin de bolsa, emisión de alarma media a los 5 segundos exactos de desvío).
\end{itemize}

\section{Escenarios Simulados y Cumplimiento de Requisitos}

De acuerdo a la especificación del proyecto, se plantearon y ejecutaron con éxito 7 escenarios de simulación que abarcan todas las exigencias temporales y lógicas del dominio. A continuación, se detalla la evidencia de ejecución para cada uno:

\subsection{Escenario 1: Funcionamiento normal}
Se verificó el arranque estable de la bomba a 50 ml/h. La simulación demuestra que la orden inicial recibida en el tiempo $t=2.00s$ se traduce en un caudal físico de infusión a los $t=4.00s$, cumpliendo la restricción temporal de arranque en menos de 3 segundos sin emitir alarmas innecesarias.
\begin{verbatim}
--- Eventos Lógicos (Auditoría) ---
Tiempo: 02.00s | Evento: NUEVA_ORDEN
Tiempo: 02.57s | Evento: CONFIRMACION_ENFERMERO
...
Tiempo | Caudal Indicado | Caudal Real Sensado
 t=00.00s |   00.00 ml/h    |    00.00 ml/h (OK)
 t=02.00s |   50.00 ml/h    |    00.00 ml/h (DESVÍO)
 t=04.00s |   50.00 ml/h    |    50.00 ml/h (OK)
 t=39.00s |   50.00 ml/h    |    50.00 ml/h (OK)
\end{verbatim}

\subsection{Escenario 2: Cambio de orden médica}
Demuestra la capacidad del sistema para transicionar de un caudal objetivo ($50~ml/h$) a uno nuevo ($80~ml/h$) en pleno funcionamiento. La nueva orden se recibe en $t=20.00s$ y se estabiliza a los $t=22.00s$.
\begin{verbatim}
Tiempo: 20.00s | Evento: NUEVA_ORDEN
...
Tiempo | Caudal Indicado | Caudal Real Sensado
 t=19.00s |   50.00 ml/h    |    50.00 ml/h (OK)
 t=20.00s |   80.00 ml/h    |    50.00 ml/h (TRANSICIÓN/DESVÍO)
 t=22.00s |   80.00 ml/h    |    80.00 ml/h (OK)
\end{verbatim}

\subsection{Escenario 3: Detención}
Valida la propiedad de seguridad al recibir una orden de infusión nula. Al recibirse la orden de $0~ml/h$ en $t=20.00s$, el controlador impone la detención, logrando el apagado físico de la bomba a los $t=22.00s$.
\begin{verbatim}
Tiempo: 20.00s | Evento: DETENCION_MEDICA
...
Tiempo | Caudal Indicado | Caudal Real Sensado
 t=19.00s |   50.00 ml/h    |    50.00 ml/h (OK)
 t=20.00s |   00.00 ml/h    |    50.00 ml/h (TRANSICIÓN/DESVÍO)
 t=22.00s |   00.00 ml/h    |    00.00 ml/h (OK)
\end{verbatim}

\subsection{Escenario 4: Desvío leve tolerado}
Este escenario inyecta una perturbación transitoria en el sensor de flujo. El caudal divergió del objetivo durante $4.0s$. Al ser estrictamente menor al umbral crítico definido por la especificación ($>5s$), el controlador asimiló la perturbación sin emitir ninguna señal acústica o visual, probando la solidez del margen de tolerancia.
\begin{verbatim}
Objetivo: Tolerancia a fallas leves...
Tiempo | Caudal Indicado | Caudal Real Sensado | Segundos de Desvío
 t=20.00s |   50.00 ml/h    |    50.00 ml/h        | 0.0 s (FALLA INYECTADA)
 t=24.00s |   50.00 ml/h    |    60.00 ml/h        | 4.0 s (OK)
 t=25.00s |   50.00 ml/h    |    50.00 ml/h        | 0.0 s (OK)
\end{verbatim}

\subsection{Escenario 5: Desvío mayor (Alarma Media)}
Ante una falla severa sostenida, la traza de simulación demuestra que exactamente al acumularse 5 segundos de discrepancia persistente (entre $t=20.00s$ y $t=25.00s$), el sistema reacciona emitiendo la señal \texttt{ALARMA\_MEDIA}, cumpliendo la restricción impuesta sobre el monitoreo continuo.
\begin{verbatim}
Tiempo: 24.11s | Evento: ALARMA_MEDIA
...
Tiempo | Caudal Indicado | Caudal Real Sensado | Segundos de Desvío
 t=20.00s |   50.00 ml/h    |    50.00 ml/h        | 0.0 s (FALLA INYECTADA)
 t=24.00s |   50.00 ml/h    |    65.00 ml/h        | 4.0 s (FALLA INYECTADA)
 t=25.00s |   50.00 ml/h    |    65.00 ml/h        | 5.0 s (ALARMA MEDIA ACTIVA)
\end{verbatim}

\subsection{Escenario 6: Fin de Bolsa y ventana de gracia}
Comprueba que, al detectar el vaciado inminente en el instante $t=20.00s$, el controlador emite una alerta preventiva e inicia el período de gracia máximo de $60~segundos$. Finalizada la ventana de gracia temporal en $t=80.00s$, el actuador detiene por completo la infusión por seguridad.
\begin{verbatim}
Tiempo: 20.00s | Evento: ALARMA_BAJA
Tiempo: 20.00s | Evento: FIN_BOLSA_DETECTADO
Tiempo: 79.11s | Evento: DETENCION_MEDICA
...
Tiempo | Caudal Real | Segundos desde Fin de Bolsa | Estado esperado
 t=20.00s |  50.00 ml/h  | 00.0 s   | ALARMA BAJA (Esperando fin gracia)
 t=80.00s |  50.00 ml/h  | 60.0 s   | DETENIDA (Pasaron 60s sin nueva bolsa)
 t=81.00s |  00.00 ml/h  | 61.0 s   | DETENIDA (Pasaron 60s sin nueva bolsa)
\end{verbatim}

\subsection{Escenario 7: Escalabilidad crítica y vivacidad}
Es el escenario más riguroso respecto a las especificaciones lógicas y temporales. Ante una falla sin corrección humana, el sistema pasa de \texttt{ALARMA\_MEDIA} ($24.11s$) a \texttt{ALARMA\_CRITICA} ($29.11s$) y detiene la infusión. Se valida con absoluta precisión matemática la restricción de temporalidad: si la alarma no es confirmada dentro de 30 segundos, comienza a repetirse (notado en $59.11s$), y luego entra en un bucle estricto cada 10 segundos ($69.11s$, $79.11s$), demostrando fehacientemente la vivacidad del diseño.
\begin{verbatim}
Tiempo: 24.11s | Evento: ALARMA_MEDIA
Tiempo: 29.11s | Evento: ALARMA_CRITICA
Tiempo: 29.11s | Evento: DETENCION_MEDICA

Historial de emisiones físicas hacia el entorno hospitalario:
 t=24.11s -> Emite Sonido/Luz: MEDIA (Primera advertencia)
 t=29.11s -> Emite Sonido/Luz: CRITICA (Debe repetirse cada 10s tras 30s)
 t=59.11s -> Emite Sonido/Luz: CRITICA (Debe repetirse cada 10s tras 30s)
 t=69.11s -> Emite Sonido/Luz: CRITICA (Debe repetirse cada 10s tras 30s)
 t=79.11s -> Emite Sonido/Luz: CRITICA (Debe repetirse cada 10s tras 30s)
\end{verbatim}